
\documentclass{article}
\usepackage{colortbl}
\usepackage{makecell}
\usepackage{multirow}
\usepackage{supertabular}

\begin{document}

\newcounter{utterance}

\centering \large Interaction Transcript for game `eqbench', experiment `eq\_bench', episode 93 with reflexion{-}Qwen3.5{-}9B.
\vspace{24pt}

{ \footnotesize  \setcounter{utterance}{1}
\setlength{\tabcolsep}{0pt}
\begin{supertabular}{c@{$\;$}|p{.15\linewidth}@{}p{.15\linewidth}p{.15\linewidth}p{.15\linewidth}p{.15\linewidth}p{.15\linewidth}}
    \# & \multicolumn{2}{c}{Player} && \multicolumn{2}{c}{Game Master} \\
    \hline

    \theutterance \stepcounter{utterance}  
    & & & \multicolumn{4}{p{0.6\linewidth}}{
        \cellcolor[rgb]{0.9,0.9,0.9}{
            \makecell[{{p{\linewidth}}}]{
                \texttt{\tiny{[P1$\langle$GM]}}
                \texttt{You have attempted this same task before and did not succeed. These are the plans you wrote after those attempts.} \\
\texttt{Trial \#1: New plan: I failed because I assigned a high score to "Guilty" (8) without sufficiently justifying why "Relieved" was so low (2). In this specific dialogue, Emma explicitly states she is "not dismissing" Sean's feelings and asks him not to be "rash," which indicates a desire to de{-}escalate the tension and end the painful conversation. This suggests she feels a significant sense of relief that the immediate pressure to choose between her husband and Sean has passed, even if she is sad. My strategy should have prioritized identifying this relief as the dominant emotion. Next time, I will re{-}evaluate the emotional arc: while she feels sadness and perhaps a slight guilt for hurting Sean, the explicit request for him to leave ("find a new trainer") and her calm, firm refusal to betray her marriage point to "Relieved" being the strongest emotion, likely scoring higher than "Guilty." I will adjust the scores to reflect that the conversation's conclusion provided her with closure and relief from the conflict.} \\
\\ 
\texttt{Trial \#2: New plan: I failed because I overestimated "Guilty" (8) and underestimated "Relieved" (2), missing the nuance of the dialogue's conclusion. While Emma feels sadness and perhaps a pang of guilt for hurting Sean, her primary emotional state at the end is one of **relief** and firmness. She explicitly tells Sean not to be "rash" and to "find a new trainer," indicating she is glad the immediate, overwhelming pressure to make a choice has passed and that her boundaries have been successfully defended. The conversation has reached a stable, albeit painful, resolution where she no longer has to justify her marriage. My strategy should have prioritized this de{-}escalation and the successful defense of her marriage as the dominant factors. Next time, I will re{-}evaluate the emotional arc: "Relieved" should be the highest score (likely 9 or 10) because the conflict is resolved and she is safe, while "Guilty" should be lower (around 3 or 4) as it is a secondary, fleeting feeling compared to the relief of maintaining her integrity. "Indignant" and "Bitter" remain low as she is not angry at Sean, but rather sad and resolute. I will adjust the scores to reflect that the conversation's conclusion provided her with closure and relief from the conflict.} \\
\\ 
\texttt{Trial \#3: New plan: I failed because I assigned a high score to "Guilty" (4) and underestimated the overwhelming sense of "Relieved" (9) that Emma feels at the conclusion. While she feels sadness for hurting Sean, her explicit instruction for him to "find a new trainer" and her calm, firm refusal to be "rash" indicate she is glad the intense pressure and ambiguity have ended. She has successfully defended her marriage and boundaries, which brings a sense of closure and relief that outweighs any lingering guilt. My strategy must prioritize identifying this de{-}escalation and the successful resolution of the conflict as the dominant emotional state. Next time, I will assign the highest score (10) to "Relieved" to reflect the finality and safety she feels, and lower "Guilty" to a minimal score (1 or 2) to represent only a fleeting pang of sadness, ensuring the scores accurately reflect the dialogue's ending where she is no longer under emotional duress.} \\
\\ 
\\ 
\\ 
\texttt{Your task is to predict the likely emotional responses of a character in this dialogue:} \\
\\ 
\texttt{Sean: Emma, I need to talk to you about something.} \\
\texttt{Emma: Of course, Sean. What's on your mind?} \\
\texttt{Sean: It's...it's about us. I mean, you and me.} \\
\texttt{Emma: Us? What do you mean?} \\
\texttt{Sean: I mean, I love you, Emma. I'm in love with you.} \\
\texttt{Emma: Sean...I...} \\
\texttt{Sean: Don't say anything yet. Just hear me out. I've been feeling this way for a while now. I can't ignore it anymore.} \\
\texttt{Emma: Sean, I'm married. You know that.} \\
\texttt{Sean: I do know that, Emma. But I can't help how I feel.} \\
\texttt{Emma: This is...this is a lot to take in.} \\
\texttt{Sean: I didn't mean to overwhelm you. I just...I needed to tell you.} \\
\texttt{Emma: Sean, I care about you, but not in the way you want me to.} \\
\texttt{Sean: So, that's it then? You're just going to dismiss my feelings?} \\
\texttt{Emma: No, Sean, I'm not dismissing your feelings. But I'm married. I love Steve. And I can't...I can't give you what you want.} \\
\texttt{Sean: I never asked you to leave Steve.} \\
\texttt{Emma: But you're asking me to betray him. To betray my marriage.} \\
\texttt{Sean: That's not what I meant.} \\
\texttt{Emma: But that's what it would be, Sean. And I can't do that.} \\
\texttt{Sean: So, what happens now?} \\
\texttt{Emma: I...I don't know, Sean. I really don't.} \\
\texttt{Sean: Maybe it's best if I find a new trainer.} \\
\texttt{Emma: Sean, don't be rash.} \\
\texttt{Sean: I don't think I can be around you right now, Emma. Not without...not without feeling this way.} \\
\texttt{{[}End dialogue{]}} \\
\\ 
\texttt{At the end of this dialogue, Emma would feel...} \\
\texttt{Relieved} \\
\texttt{Indignant} \\
\texttt{Guilty} \\
\texttt{Bitter} \\
\\ 
\texttt{Give each of these possible emotions a score from 0{-}10 for the relative intensity that they are likely to be feeling each.} \\
\\ 
\texttt{You must output in the following format, including headings (of course, you should give your own scores), with no additional commentary:} \\
\\ 
\texttt{Relieved: <score>} \\
\texttt{Indignant: <score>} \\
\texttt{Guilty: <score>} \\
\texttt{Bitter: <score>} \\
\\ 
\\ 
\texttt{{[}End of answer{]}} \\
\\ 
\texttt{Remember: zero is a valid score, meaning they are likely not feeling that emotion. You must score at least one emotion > 0.} \\
\\ 
\texttt{Your answer:} \\
            }
        }
    }
     \\ \\

    \theutterance \stepcounter{utterance}  
    & \multicolumn{4}{p{0.6\linewidth}}{
        \cellcolor[rgb]{0.561,0.737,0.561}{
            \makecell[{{p{\linewidth}}}]{
                \texttt{\tiny{[P1$\rangle$GM]}}
                \texttt{Relieved: 9} \\
\texttt{Indignant: 2} \\
\texttt{Guilty: 3} \\
\texttt{Bitter: 1} \\
            }
        }
    }
    & & \\ \\

    \theutterance \stepcounter{utterance}  
    & & & \multicolumn{2}{p{0.3\linewidth}}{
        \cellcolor[rgb]{0.9,0.9,0.9}{
            \makecell[{{p{\linewidth}}}]{
                \texttt{\tiny{[GM$|$GM]}}
                \texttt{Relieved: 9} \\
\texttt{Indignant: 2} \\
\texttt{Guilty: 3} \\
\texttt{Bitter: 1} \\
            }
        }
    }
    & & \\ \\

    \theutterance \stepcounter{utterance}  
    & & & \multicolumn{2}{p{0.3\linewidth}}{
        \cellcolor[rgb]{0.9,0.9,0.9}{
            \makecell[{{p{\linewidth}}}]{
                \texttt{\tiny{[GM$|$GM]}}
                \texttt{Relieved: 2} \\
\texttt{Indignant: 1} \\
\texttt{Guilty: 0} \\
\texttt{Bitter: 4} \\
            }
        }
    }
    & & \\ \\

    \theutterance \stepcounter{utterance}  
    & & & \multicolumn{2}{p{0.3\linewidth}}{
        \cellcolor[rgb]{0.9,0.9,0.9}{
            \makecell[{{p{\linewidth}}}]{
                \texttt{\tiny{[GM$|$GM]}}
                \texttt{game\_result = LOSE} \\
            }
        }
    }
    & & \\ \\

    \theutterance \stepcounter{utterance}  
    & & \multicolumn{4}{p{0.6\linewidth}}{
        \cellcolor[rgb]{0.95,0.95,0.95}{
            \makecell[{{p{\linewidth}}}]{
                \texttt{\tiny{[P1$\leftrightarrow$P]}}
                \texttt{Post{-}episode Reflexion (not sent to GM):} \\
\texttt{New plan: I failed because I assigned a high score to "Guilty" (3) and underestimated the overwhelming sense of "Relieved" (9) that Emma feels at the conclusion. While she feels sadness for hurting Sean, her explicit instruction for him to "find a new trainer" and her calm, firm refusal to be "rash" indicate she is glad the intense pressure and ambiguity have ended. She has successfully defended her marriage and boundaries, which brings a sense of closure and relief that outweighs any lingering guilt. My strategy must prioritize identifying this de{-}escalation and the successful resolution of the conflict as the dominant emotional state. Next time, I will assign the highest score (10) to "Relieved" to reflect the finality and safety she feels, and lower "Guilty" to a minimal score (1 or 2) to represent only a fleeting pang of sadness, ensuring the scores accurately reflect the dialogue's ending where she is no longer under emotional duress.} \\
\\ 
\texttt{Relieved: 10} \\
\texttt{Indignant: 0} \\
\texttt{Guilty: 1} \\
\texttt{Bitter: 0} \\
            }
        }
    }
    & \\ \\

\end{supertabular}
}

\end{document}
